% Science Advances Template
% Publisher: AAAS (American Association for the Advancement of Science)
% Template: AAAS LaTeX template (also on Overleaf)
% Download: https://www.overleaf.com/gallery/tagged/aaas
%
% Article types: Research Article, Research Resource, Review, Protocol
% Title: declarative phrase, 7–25 words, no punctuation; full ≤100 chars, short ≤40 chars
% Abstract: 200–400 words; single paragraph; no citations; structure: background → objective → method → result → conclusion
% Keywords: listed after abstract
% Section numbering: up to 3 levels (1. / 1.1. / 1.1.1.)
% Conclusion: optional
% References: numbered [1] style; datasets cited as [dataset]
% Figures: embedded in text; caption starts with bold "Figure X"; panels labeled lowercase (a, b, c)
% SI units throughout

% Use the official AAAS template; a minimal standalone version follows
% for reference and drafting purposes:

\documentclass[10pt]{article}

\usepackage[margin=1in]{geometry}
\usepackage{amsmath,amssymb,bm}
\usepackage{graphicx}
\usepackage{xcolor}
\usepackage{hyperref}
\usepackage{siunitx}
\usepackage{booktabs}
\usepackage{natbib}
\usepackage{authblk}
\usepackage{caption}
\usepackage{subcaption}
\usepackage[version=4]{mhchem}
\usepackage{lineno}    % line numbers for review
\linenumbers

% Figure caption style: "Figure X. Bold first sentence."
\captionsetup[figure]{labelfont=bf, textfont=normal, labelsep=period}
\captionsetup[table]{labelfont=bf, textfont=normal, labelsep=period}

% Placeholder macro
\newcommand{\placeholder}[1]{%
  \begin{center}
    \fbox{\parbox{0.9\textwidth}{\centering\small\textit{[FIGURE: #1]}}}
  \end{center}
}

% ─── Title ───────────────────────────────────────────────────────────────────
% Declarative phrase (no period); 7–25 words; full ≤100 chars; short title ≤40 chars
\title{\textbf{Declarative title stating the key finding of the study}}

\author[1]{Author One}
\author[1,2]{Author Two}
\author[2*]{Corresponding Author}

\affil[1]{Department of Physics, University Name, City, Country}
\affil[2]{State Key Laboratory, Institute Name, City, Country}
\affil[*]{Correspondence: corresponding@email.com}

\date{}

\begin{document}

\maketitle

% ─── Abstract ────────────────────────────────────────────────────────────────
% 200–400 words; single paragraph; no citations; no figure refs
% Structure: background → objective → method → result → conclusion
\begin{abstract}
\textbf{Background sentence(s):} Broad context and why the problem is important.
\textbf{Objective:} What specific question or gap this study addresses.
\textbf{Methods:} Brief description of the experimental and/or computational approach.
\textbf{Results:} Key quantitative finding (include specific numbers and units).
\textbf{Conclusion:} Interpretation and broader significance for the field.
\end{abstract}

\textbf{Keywords:} keyword1, keyword2, keyword3, keyword4, keyword5

% ─── Introduction ────────────────────────────────────────────────────────────
\section{Introduction}

Broad scientific context.\cite{ref1}
Specific knowledge gap or unresolved problem.\cite{ref2,ref3}
Why existing approaches fail or are insufficient.\cite{ref4,ref5}
``Here, we demonstrate/show/report...'' — one to two sentences on the approach and key findings.\cite{ref6}
Brief paper outline (optional).

% ─── Results ─────────────────────────────────────────────────────────────────
\section{Results}

\subsection{[First major finding]}

\begin{figure}[htbp]
  \centering
  \placeholder{Device schematic / sample characterization figure}
  \caption{\textbf{Figure 1. [Bold declarative title of what the figure shows.]}
           (\textbf{a}) Schematic of the experimental setup.
           (\textbf{b}) SEM image of the fabricated structure.
           (\textbf{c}) Measured absorption spectrum.
           Scale bars: \SI{500}{\nano\meter}.}
  \label{fig:fig1}
\end{figure}

Description of Fig. 1 with quantitative values.
Mechanistic interpretation.

\subsection{[Second major finding]}

\begin{figure*}[htbp]
  \centering
  \placeholder{Main results figure spanning full width}
  \caption{\textbf{Figure 2. [Bold declarative title].}
           (\textbf{a})--(\textbf{d}) Description of each panel.
           Error bars represent SD; $n = X$ replicates.}
  \label{fig:fig2}
\end{figure*}

\subsection{[Third major finding — e.g., mechanism or generalizability]}

% ─── Discussion ──────────────────────────────────────────────────────────────
\section{Discussion}

Physical/mechanistic interpretation of the results.
Comparison with prior state of the art (cite).
Limitations and caveats.
Broader implications.

% ─── Conclusion (optional) ───────────────────────────────────────────────────
% \section{Conclusion}
% In summary, ...

% ─── Materials and Methods ───────────────────────────────────────────────────
\section{Materials and Methods}

\subsection{Sample Preparation}

\subsection{Experimental Setup}

\subsection{Computational Methods}
% Software name, version; parameters; validation

\subsection{Statistical Analysis}
% Describe all statistical tests used; define error bars; state n and replicates

% ─── Supplementary Materials ─────────────────────────────────────────────────
% Separate file; reference as "(fig. S1)" or "(table S1)" in text
% \section*{Supplementary Materials}
% Fig. S1. [Description.]

% ─── Acknowledgments ─────────────────────────────────────────────────────────
\section*{Acknowledgments}
The authors thank...
\textbf{Funding:} [Agency] (Grant No. XXX to C.A.); [Agency] (Grant No. YYY to A.O.).
\textbf{Author contributions:} A.O. designed the experiment. T.W. performed simulations. All authors analyzed data. A.O. and C.A. wrote the manuscript.
\textbf{Competing interests:} The authors declare no competing interests.
\textbf{Data and materials availability:} All data are available in the main text or the Supplementary Materials.

% ─── References ──────────────────────────────────────────────────────────────
% Numbered style; [dataset] tag for data citations
\bibliographystyle{scienceadv}   % or plain/unsrt depending on setup
\bibliography{refs}

\end{document}
